% Shared preamble: packages and typography for the thesis AND the playground,
% so whatever works in the playground works in the thesis.
% Page geometry is a placeholder until the PG editorial guidelines are in.
%
% A document may switch author notes off BEFORE loading this file:
%   \newif\ifworkingdraft \workingdraftfalse
% Otherwise it is a working draft (notes and the "wersja robocza" footer on).
% (\csname keeps a literal \ifworkingdraft out of the skipped branch, which
% TeX would otherwise count as an unmatched conditional.)
\ifcsname ifworkingdraft\endcsname\else
  \expandafter\newif\csname ifworkingdraft\endcsname
  \workingdrafttrue
\fi

% --- Language and fonts -------------------------------------------------------
\usepackage[english, main=polish]{babel}
\usepackage{amsmath}
\usepackage{fontspec}
\usepackage{unicode-math}
\setmainfont{TeX Gyre Termes}
\setsansfont{TeX Gyre Heros}
\setmonofont{DejaVu Sans Mono}[Scale=MatchLowercase]
\setmathfont{TeX Gyre Termes Math}
\usepackage[autostyle]{csquotes}   % \enquote{...} -> „...”
\usepackage{microtype}
\usepackage{indentfirst}           % Polish: indent the first paragraph too

% One-letter words (w, z, i, a, o, u) never end a line — no manual "~".
\usepackage{luavlna}
\singlechars{polish}{AaIiOoUuWwZz}

% --- Page (placeholder until the PG guidelines) -------------------------------
\usepackage[a4paper, margin=2.5cm, marginparwidth=2.1cm]{geometry}
\usepackage[onehalfspacing]{setspace}
\usepackage{xcolor}

\usepackage{fancyhdr}
\fancypagestyle{plain}{%
  \fancyhf{}%
  \fancyfoot[C]{\thepage}%
  \ifworkingdraft\fancyfoot[R]{\scriptsize\color{gray}wersja robocza, \today}\fi
  \renewcommand{\headrulewidth}{0pt}%
}
\pagestyle{plain}

% --- Headings (Polish convention: a full stop after the number) --------------
\usepackage{titlesec}
\titleformat{\chapter}[hang]{\normalfont\huge\bfseries}{\thechapter.}{0.6em}{}
\titlespacing*{\chapter}{0pt}{0pt}{1.5em}
\titleformat{\section}{\normalfont\Large\bfseries}{\thesection.}{0.6em}{}
\titleformat{\subsection}{\normalfont\large\bfseries}{\thesubsection.}{0.6em}{}

% --- Floats, tables, numbers ---------------------------------------------------
\usepackage{graphicx}
\graphicspath{{figures/}}
\usepackage{booktabs}
\usepackage{tabularx}
\usepackage{caption}
\captionsetup{font=small, labelfont=bf, labelsep=period}
\captionsetup[table]{position=above, skip=6pt}  % table caption ABOVE
\captionsetup[figure]{position=below}           % figure caption BELOW
\usepackage{siunitx}
\sisetup{output-decimal-marker={,}, group-separator={\,}, group-minimum-digits=4}
\usepackage{enumitem}
\usepackage{tikz}
\usepackage{pdfpages}

% --- Code ----------------------------------------------------------------------
\usepackage{listings}
\lstset{
  basicstyle=\ttfamily\footnotesize,
  keywordstyle=\color{blue!60!black}\bfseries,
  commentstyle=\color{gray}\itshape,
  stringstyle=\color{red!60!black},
  numbers=left, numberstyle=\tiny\color{gray},
  frame=tb, breaklines=true, showstringspaces=false,
  columns=fullflexible, keepspaces=true, captionpos=t,
}
\renewcommand{\lstlistingname}{Listing}
\renewcommand{\lstlistlistingname}{Spis listingów}
\newcommand{\code}{\lstinline[basicstyle=\ttfamily\small]}

% --- Bibliography (style is a placeholder until the supervisor decides) -------
\usepackage[backend=biber, style=numeric-comp, sorting=none]{biblatex}

% --- Acronyms: variant B ------------------------------------------------------
% \ac{x} always prints the short form. Write the first use out by hand in the
% right grammatical case: "...w temperaturze pokojowej (\ac{tp})...".
\usepackage{acro}
\acsetup{first-style=short, make-links=true}

% --- Author notes in the margin, hidden in the final version ----------------
\ifworkingdraft
  \usepackage[textsize=scriptsize]{todonotes}
\else
  \usepackage[disable]{todonotes}
\fi
\newcommand{\marcin}[1]{\todo[color=blue!25, author=Marcin]{#1}}
\newcommand{\oskar}[1]{\todo[color=green!25, author=Oskar]{#1}}
\newcommand{\kacper}[1]{\todo[color=orange!30, author=Kacper]{#1}}
\newcommand{\supervisor}[1]{\todo[color=red!25, author=Promotor]{#1}}

% --- Links (last) -------------------------------------------------------------
% hypertexnames=false: the title page and the table of contents are both
% "page 1", which would otherwise give duplicate link targets.
\usepackage[unicode, hidelinks, hypertexnames=false]{hyperref}
